
\documentclass[11pt,a4paper]{article}

\usepackage[utf8]{inputenc}
\usepackage[T1]{fontenc}
\usepackage[french]{babel}
\usepackage{lmodern}
\usepackage{microtype}
\usepackage[a4paper,margin=1.65cm,headheight=15pt]{geometry}
\usepackage{amsmath,amssymb,mathtools}
\usepackage{array,tabularx,multicol}
\usepackage{enumitem}
\usepackage[most]{tcolorbox}
\usepackage{xcolor}
\usepackage{tikz}
\usetikzlibrary{arrows.meta,calc}
\usepackage{pdflscape}
\usepackage{fancyhdr}
\usepackage{fancyvrb}
\usepackage{listings}
\usepackage{hyperref}

\definecolor{fondpage}{RGB}{247,248,250}
\definecolor{encre}{RGB}{31,35,40}
\definecolor{bleu}{RGB}{40,91,135}
\definecolor{vert}{RGB}{55,125,92}
\definecolor{orange}{RGB}{178,105,42}
\definecolor{violet}{RGB}{101,77,145}
\definecolor{rouge}{RGB}{170,72,72}
\definecolor{gris}{RGB}{86,91,99}
\definecolor{bleufonce}{RGB}{28,55,120}
\definecolor{bleuclair}{RGB}{238,243,255}
\definecolor{grisclair}{RGB}{245,245,245}
\definecolor{tableline}{RGB}{45,45,45}
\definecolor{softgray}{RGB}{246,247,249}
\definecolor{deepblue}{RGB}{25,62,115}
\definecolor{softblue}{RGB}{235,241,250}
\definecolor{accent}{RGB}{20,82,130}

\hypersetup{colorlinks=true,linkcolor=bleufonce,urlcolor=bleufonce,pdftitle={Parcours de révision -- Géométrie dans l'espace -- Terminale spécialité},pdfauthor={}}

\setlength{\parindent}{0pt}
\setlength{\parskip}{0.25em}
\renewcommand{\arraystretch}{1.12}
\setlength{\tabcolsep}{3pt}
\setlist[itemize]{leftmargin=1.25em,itemsep=0.15em,topsep=0.2em}
\setlist[enumerate,1]{label=\textbf{\arabic*.}, leftmargin=1.05cm, itemsep=0.3em, topsep=0.2em}
\setlist[enumerate,2]{label=\textbf{\alph*.}, leftmargin=0.85cm, itemsep=0.2em, topsep=0.15em}

\pagestyle{fancy}
\fancyhf{}
\lhead{}
\rhead{}
\cfoot{\thepage}

\newcommand{\R}{\mathbb{R}}
\newcommand{\N}{\mathbb{N}}
\newcommand{\vect}[1]{\overrightarrow{#1}}
\newcommand{\vu}{\vec u}
\newcommand{\vv}{\vec v}
\newcommand{\vw}{\vec w}
\newcommand{\vn}{\vec n}
\newcommand{\vd}{\vec d}
\newcommand{\norme}[1]{\left\lVert #1\right\rVert}
\newcommand{\ps}{\cdot}
\newcommand{\e}{\mathrm e}
\newcommand{\origine}[1]{ {\scriptsize\textcolor{bleufonce}{(site #1)}}}

\newcounter{question}
\newcounter{reponse}
\newenvironment{blocquestions}[1]{%
  \begin{tcolorbox}[
    enhanced,
    breakable,
    colback=bleuclair,
    colframe=bleufonce,
    coltitle=white,
    colbacktitle=bleufonce,
    title={\bfseries #1},
    fonttitle=\bfseries,
    boxrule=0.6pt,
    arc=2mm,
    left=2mm,
    right=2mm,
    top=2.5mm,
    bottom=1.5mm,
    toptitle=1.2mm,
    bottomtitle=1.2mm,
    lefttitle=2mm,
    righttitle=2mm,
    before upper=\vspace{0.7mm},
    before skip=3.5mm,
    after skip=3mm, fontupper=\large
  ]
  \large
  \begin{enumerate}[leftmargin=*,label=\textbf{\arabic*.},itemsep=0.85em,topsep=0.5em]
  \setcounter{enumi}{\value{question}}
}{%
  \setcounter{question}{\value{enumi}}
  \end{enumerate}
  \end{tcolorbox}
}
\newenvironment{blocreponses}[1]{%
  \begin{tcolorbox}[
    enhanced,
    breakable,
    colback=bleuclair,
    colframe=bleufonce,
    coltitle=white,
    colbacktitle=bleufonce,
    title={\bfseries #1},
    fonttitle=\bfseries,
    boxrule=0.6pt,
    arc=2mm,
    left=2mm,
    right=2mm,
    top=2.5mm,
    bottom=1.5mm,
    toptitle=1.2mm,
    bottomtitle=1.2mm,
    lefttitle=2mm,
    righttitle=2mm,
    before upper=\vspace{0.7mm},
    before skip=3.5mm,
    after skip=3mm, fontupper=\large
  ]
  \large
  \begin{enumerate}[leftmargin=*,label=\textbf{\arabic*.},itemsep=0.45em,topsep=0.5em]
  \setcounter{enumi}{\value{reponse}}
}{%
  \setcounter{reponse}{\value{enumi}}
  \end{enumerate}
  \end{tcolorbox}
}

\newtcolorbox{correctionbox}[1]{enhanced,breakable,colback=grisclair,colframe=black!55,boxrule=0.55pt,arc=2mm,left=2mm,right=2mm,top=1.5mm,bottom=1.5mm,title=\textbf{#1},coltitle=black,colbacktitle=black!12,before skip=2mm,after skip=2mm}
\newcommand{\titreSujet}[2]{%
  \subsection*{#1}
  \addcontentsline{toc}{subsection}{#1}
  \textit{Référence : #2}\par\medskip\hrule\medskip
}

% Styles de la carte mentale de cours
\setlist[itemize]{
    leftmargin=1.1em,
    itemsep=0.10em,
    topsep=0.15em,
    parsep=0pt,
    partopsep=0pt,
    label=\textcolor{black!65}{\raisebox{0.15ex}{\tiny$\blacktriangleright$}}
}
\tcbset{
    enhanced,
    arc=2.2mm,
    outer arc=2.2mm,
    boxrule=0.45pt,
    left=1.55mm,
    right=1.55mm,
    top=1.65mm,
    bottom=1.65mm,
    boxsep=0.8mm,
    before skip=1.6mm,
    after skip=1.6mm,
    fonttitle=\bfseries\small,
    coltitle=encre,
    before upper=\raggedright\fontsize{9.05}{10.9}\selectfont,
}
\newtcolorbox{fichebox}[3][]{
    colback=white,
    colframe=#2!72!black,
    colbacktitle=#2!15,
    coltitle=#2!55!black,
    borderline west={1mm}{0pt}{#2!78!black},
    title={#3},
    drop fuzzy shadow=black!4,
    #1
}
\newcommand{\tagcol}[2]{%
\begin{tcolorbox}[
    enhanced,
    colback=#1!11,
    colframe=#1!45!black,
    boxrule=0.42pt,
    arc=2.2mm,
    left=1mm,right=1mm,top=0.55mm,bottom=0.55mm,
    before skip=0mm,
    after skip=1mm,
    fontupper=\bfseries\small,
    colupper=#1!55!black
]
\centering #2
\end{tcolorbox}}
\newcommand{\titrepage}{%
\begin{tcolorbox}[
    colback=white,
    colframe=bleu!70!black,
    boxrule=0.65pt,
    arc=3mm,
    left=2.5mm,right=2.5mm,top=1.2mm,bottom=1.2mm,
    before skip=0mm,
    after skip=1.6mm,
    drop fuzzy shadow=black!8
]
{\Large\bfseries Parcours de révision -- Géométrie dans l'espace}\hfill
{\large\bfseries Terminale spécialité}

\vspace{0.25mm}
{\small Repérage, droites, plans, produit scalaire, orthogonalité et distances}
\end{tcolorbox}}

\tikzset{tabouter/.style={draw=tableline, line width=0.85pt},tabline/.style={draw=tableline, line width=0.55pt},forbidden/.style={draw=tableline, line width=0.95pt},vararrow/.style={-{Latex[length=2.7mm,width=2.0mm]}, line width=0.85pt, draw=accent},tabtext/.style={font=\small},tabmath/.style={font=\small},labelcell/.style={font=\small\bfseries, text=deepblue},pointvalue/.style={font=\small\bfseries},endvalue/.style={font=\small}}
\lstset{language=Python,basicstyle=\ttfamily\small,numbers=left,numberstyle=\tiny,frame=single,showstringspaces=false,columns=fullflexible,keepspaces=true}


\begin{document}
\begin{center}
{\LARGE\bfseries Parcours de révision -- Géométrie dans l'espace}\\[1mm]

\end{center}
\vspace{2mm}
\tableofcontents
\clearpage

\section*{Partie 1 -- Corrections des questions flashs}
\addcontentsline{toc}{section}{Partie 1 -- Corrections des questions flashs}
\setcounter{reponse}{0}\begin{blocreponses}{Réponses -- A. Repères, vecteurs, droites et plans}
  \item On a $\overrightarrow{AB}=(x_B-x_A\,;\,y_B-y_A\,;\,z_B-z_A)$.
  \item On fait la moyenne des coordonnées : $I\left(\dfrac{x_A+x_B}{2};\dfrac{y_A+y_B}{2};\dfrac{z_A+z_B}{2}\right)$.
  \item On peut montrer que les vecteurs $\overrightarrow{AB}$ et $\overrightarrow{AC}$ sont colinéaires.
  \item Cela signifie que l'un est un multiple de l'autre.
  \item C'est un vecteur non nul colinéaire à cette droite.
  \item C'est un système donnant $x$, $y$ et $z$ en fonction d'un paramètre réel.
  \item Il permet de parcourir les différents points de la droite.
  \item On cherche s'il existe une valeur du paramètre qui donne ses coordonnées.
  \item C'est un vecteur perpendiculaire à tous les vecteurs du plan.
  \item Une équation cartésienne du plan est de la forme $ax+by+cz+d=0$.
  \item On remplace $x$, $y$ et $z$ par les coordonnées du point et on vérifie si l'égalité est vraie.
  \item On résout le système formé par les deux équations ; on obtient souvent une représentation paramétrique de l'intersection.
  \item Deux plans sont parallèles lorsque leurs vecteurs normaux sont colinéaires.
  \item Une droite est parallèle à un plan lorsque son vecteur directeur est orthogonal à un vecteur normal du plan.
  \item Cela signifie qu'ils appartiennent à un même plan.
  \item On peut montrer qu'il n'existe pas de réels $a$ et $b$ tels que $\overrightarrow{AD}=a\overrightarrow{AB}+b\overrightarrow{AC}$.
\end{blocreponses}

\begin{blocreponses}{Réponses -- B. Distances, produit scalaire et normes}
  \item $V=\dfrac13\times \text{aire de la base}\times \text{hauteur}$.
  \item $\vec u\cdot\vec v=xx'+yy'+zz'$.
  \item $\vec u\cdot\vec u=\|\vec u\|^2$.
  \item Ils sont orthogonaux lorsque leur produit scalaire est nul.
  \item $\vec u\cdot\vec v=\|\vec u\|\,\|\vec v\|\cos(\theta)$.
  \item $\|(x,y,z)\|=\sqrt{x^2+y^2+z^2}$.
  \item On calcule la norme du vecteur $\overrightarrow{AB}$ : $AB=\|\overrightarrow{AB}\|$.
  \item Une droite est orthogonale à un plan lorsque son vecteur directeur est colinéaire à un vecteur normal du plan.
  \item C'est le point du plan le plus proche du point donné.
  \item La longueur $AH$ représente la distance du point $A$ au plan $(P)$.
  \item Un vecteur directeur est $(2;-1;5)$.
  \item Un vecteur normal est $(2;-1;3)$.
  \item On peut conclure que les deux vecteurs sont orthogonaux.
\end{blocreponses}

\begin{blocreponses}{Réponses -- C. Définitions et rédaction}
  \item Ce sont deux droites contenues dans un même plan.
  \item C'est une équation de la forme $ax+by+cz+d=0$.
  \item C'est sa longueur.
  \item Un vecteur n'a pas d'équation. On écrit plutôt : « Le vecteur $\vec u$ a pour coordonnées $(1;2;3)$. »
\end{blocreponses}

\begin{blocreponses}{Réponses -- D. Calculs directs}
  \item $\overrightarrow{AB}=(4-1\,;\,0-2\,;\,5-3)=(3;-2;2)$.
  \item $(1;0;2)\cdot(3;-1;1)=1\times3+0\times(-1)+2\times1=5$.
  \item Il suffit de montrer que $\overrightarrow{AB}\cdot\overrightarrow{AC}=0$.
  \item Dans l'espace, c'est la sphère de diamètre $[AB]$ ; dans un plan, c'est le cercle de diamètre $[AB]$.
  \item Dans le plan, on écrit $ax+by+c=0$ avec $(a;b)$ vecteur normal, puis on utilise les coordonnées du point pour déterminer $c$. Dans l'espace, un vecteur normal détermine plutôt un plan.
\end{blocreponses}
\clearpage
\section*{Partie 2 -- Corrigés des sujets de bac}
\addcontentsline{toc}{section}{Partie 2 -- Corrigés des sujets de bac}
\titreSujet{Sujet 1 -- Vrai / faux dans l'espace}{Amérique du Nord, 21 mai 2025, jour 1, exercice 3}

\begin{correctionbox}{Affirmation 1}
Un vecteur directeur de la droite $(d)$ est $\vec u(-2;0;-6)$. L'axe des ordonnées a pour vecteur directeur $\vec j(0;1;0)$ et passe par $O(0;0;0)$.
Le point $M(3;-1;2)$ appartient à $(d)$, donc $\vect{OM}(3;-1;2)$.
On calcule le produit mixte :
\[
\det(\vec u,\vec j,\vect{OM})=14\neq 0.
\]
Les deux droites ne sont donc pas coplanaires. L'affirmation est \textbf{vraie}.
\end{correctionbox}

\begin{correctionbox}{Affirmation 2}
Un plan orthogonal à la droite $(d)$ admet pour vecteur normal un vecteur directeur de $(d)$, par exemple $\vec n(1;0;3)$.
Le plan passant par $A(3;-3;-2)$ a donc une équation de la forme
\[
(x-3)+3(z+2)=0,
\]
c'est-à-dire
\[
x+3z+3=0.
\]
L'affirmation est \textbf{vraie}.
\end{correctionbox}

\begin{correctionbox}{Affirmation 3}
Le point $C$ est le point de $(d)$ d'abscisse $2$. On résout $3-2t=2$, donc $t=\frac12$, puis
\[
C(2;-1;-1).
\]
On a
\[
\vect{AB}(2;-1;1),\qquad \vect{AC}(-1;2;1).
\]
Ainsi
\[
\vect{AB}\cdot\vect{AC}=2\times(-1)+(-1)\times2+1\times1=-3.
\]
De plus $AB=AC=\sqrt6$, donc
\[
\cos(\widehat{BAC})=\dfrac{-3}{\sqrt6\sqrt6}=-\dfrac12.
\]
Donc $\widehat{BAC}=\dfrac{2\pi}{3}$ et non $\dfrac{\pi}{6}$. L'affirmation est \textbf{fausse}.
\end{correctionbox}

\begin{correctionbox}{Affirmation 4}
Le point $H$ est le projeté orthogonal de $B$ sur le plan $\mathcal P$, donc $BH$ est la distance du point $B(5;-4;-1)$ au plan $x+3z-7=0$.
On utilise la formule de distance d'un point à un plan :
\[
BH=\dfrac{|5+3(-1)-7|}{\sqrt{1^2+0^2+3^2}}
=\dfrac{5}{\sqrt{10}}=\dfrac{\sqrt{10}}2.
\]
L'affirmation est \textbf{vraie}.
\end{correctionbox}

\newpage
\titreSujet{Sujet 2 -- Plan, projeté orthogonal et tétraèdre}{Polynésie, 18 juin 2025, sujet 2, exercice 3}

\begin{correctionbox}{1. Plan déterminé par $A$, $B$ et $C$}
On calcule
\[
\vect{AB}(-2;1;5),\qquad \vect{AC}(-1;-2;0).
\]
Ces deux vecteurs ne sont pas colinéaires, par exemple les rapports $\frac{-2}{-1}$ et $\frac{1}{-2}$ ne sont pas égaux. Les points $A$, $B$ et $C$ ne sont donc pas alignés et déterminent un plan.
\end{correctionbox}

\begin{correctionbox}{2. Triangle rectangle}
On calcule
\[
\vect{AB}\cdot\vect{AC}=(-2)(-1)+1(-2)+5\times0=0.
\]
Les vecteurs $\vect{AB}$ et $\vect{AC}$ sont orthogonaux, donc le triangle $ABC$ est rectangle en $A$.
\end{correctionbox}

\begin{correctionbox}{3. Droite orthogonale et équation du plan}
\begin{enumerate}
\item Le vecteur directeur de $\Delta$ est $\vec u(2;-1;1)$. On vérifie :
\[
\vec u\cdot\vect{AB}=2(-2)+(-1)\times1+1\times5=0,
\]
\[
\vec u\cdot\vect{AC}=2(-1)+(-1)(-2)+1\times0=0.
\]
Le vecteur $\vec u$ est orthogonal à deux vecteurs non colinéaires du plan $(ABC)$, donc $\Delta$ est orthogonale au plan $(ABC)$.

\item Le vecteur $\vec u(2;-1;1)$ est normal au plan. Une équation du plan est donc
\[
2x-y+z+d=0.
\]
Comme $A(1;3;0)$ appartient au plan :
\[
2-3+0+d=0,
\]
donc $d=1$. Ainsi
\[
(ABC):\quad 2x-y+z+1=0.
\]

\item La droite $\Delta$ passe par $D(-2;2;1)$ et a pour vecteur directeur $(2;-1;1)$, donc une représentation paramétrique est
\[
\begin{cases}
x=-2+2t,\\
y=2-t,\\
z=1+t,
\end{cases}\qquad t\in\mathbb R.
\]
\end{enumerate}
\end{correctionbox}

\begin{correctionbox}{4. Projeté orthogonal}
On vérifie d'abord que $H$ appartient au plan :
\[
2\left(-\dfrac23\right)-\dfrac43+\dfrac53+1=0.
\]
De plus
\[
\vect{DH}\left(\dfrac43;-\dfrac23;\dfrac23\right)=\dfrac23(2;-1;1),
\]
donc $\vect{DH}$ est normal au plan $(ABC)$. Ainsi $H$ est bien le projeté orthogonal de $D$ sur le plan $(ABC)$.
\end{correctionbox}

\begin{correctionbox}{5. Volume du tétraèdre}
\begin{enumerate}
\item On a
\[
DH=\sqrt{\left(\dfrac43\right)^2+\left(-\dfrac23\right)^2+\left(\dfrac23\right)^2}
=\sqrt{\dfrac{24}{9}}=\dfrac{2\sqrt6}{3}.
\]

\item Comme $ABC$ est rectangle en $A$,
\[
\mathcal A_{ABC}=\dfrac12 AB\times AC.
\]
Or
\[
AB=\sqrt{30},\qquad AC=\sqrt5,
\]
donc
\[
\mathcal A_{ABC}=\dfrac12\sqrt{150}=\dfrac{5\sqrt6}{2}.
\]
La hauteur relative à la base $ABC$ est $DH$. Donc
\[
V=\dfrac13\times \dfrac{5\sqrt6}{2}\times \dfrac{2\sqrt6}{3}=\dfrac{10}{3}.
\]
\end{enumerate}
\end{correctionbox}

\begin{correctionbox}{6. Position relative de $d$ et du plan}
On remplace les coordonnées paramétriques de $d$ dans l'équation du plan :
\[
2(1-2k)-(-3k)+(1+k)+1=4.
\]
Cette expression ne dépend pas de $k$ et n'est jamais égale à $0$. La droite $d$ ne coupe donc pas le plan $(ABC)$. Elle est parallèle au plan, et même strictement parallèle.
\end{correctionbox}

\newpage
\titreSujet{Sujet 3 -- Drones et plan dans l'espace}{Polynésie, 20 juin 2024, sujet 2, exercice 4}

\begin{correctionbox}{1. Non-alignement de $A$, $B$ et $C$}
On calcule
\[
\vect{AB}(5;-1;-13),\qquad \vect{AC}(2;-2;-10).
\]
Ces deux vecteurs ne sont pas colinéaires. Les points $A$, $B$ et $C$ ne sont donc pas alignés.
\end{correctionbox}

\begin{correctionbox}{2. Plan $\mathcal P$}
\begin{enumerate}
\item Avec $\vec n(2;-3;1)$, on a
\[
\vec n\cdot\vect{AB}=2\times5+(-3)(-1)+1(-13)=0,
\]
\[
\vec n\cdot\vect{AC}=2\times2+(-3)(-2)+1(-10)=0.
\]
Le vecteur $\vec n$ est orthogonal à deux vecteurs non colinéaires du plan, donc il est normal au plan $(ABC)$.

\item Une équation du plan est $2x-3y+z+d=0$. Comme $A(-1;-1;17)$ appartient au plan :
\[
-2+3+17+d=0,
\]
donc $d=-18$. Ainsi
\[
\mathcal P:\quad 2x-3y+z-18=0.
\]
\end{enumerate}
\end{correctionbox}

\begin{correctionbox}{3. Quatrième drone}
\begin{enumerate}
\item Un vecteur directeur de $d$ est
\[
\vec u(3;1;4).
\]
\item On cherche l'intersection avec $\mathcal P$. On remplace :
\[
2(3t+2)-3(t+5)+(4t+1)-18=0.
\]
Cela donne
\[
7t-28=0,
\]
donc $t=4$. Ainsi
\[
E(3\times4+2;4+5;4\times4+1)=(14;9;17).
\]
\end{enumerate}
\end{correctionbox}

\begin{correctionbox}{4. Distance du point $D$ au plan}
Le point $F(6;-1;3)$ est l'intersection de la droite perpendiculaire au plan passant par $D$ avec le plan. La distance du point $D$ au plan est donc $DF$.
On calcule
\[
\vect{DF}(4;-6;2),
\]
donc
\[
DF=\sqrt{4^2+(-6)^2+2^2}=\sqrt{56}=2\sqrt{14}.
\]
La distance de $D$ au plan $\mathcal P$ vaut donc $2\sqrt{14}$ centaines de mètres.
\end{correctionbox}

\begin{correctionbox}{5. Arrivée à temps}
La distance minimale à parcourir pour aller de $D$ à un point du plan $\mathcal P$ est la distance de $D$ au plan, soit
\[
2\sqrt{14}\text{ centaines de mètres}=200\sqrt{14}\text{ mètres}\approx 748{,}3\text{ m}.
\]
En $40$ secondes à $18{,}6$ m.s$^{-1}$, le drone parcourt au maximum
\[
18{,}6\times40=744\text{ m}.
\]
Comme $744<748{,}3$, le nouveau drone ne peut pas arriver à temps.
\end{correctionbox}
\end{document}
